\documentclass[a4paper]{article}

\usepackage{graphicx}
\usepackage{amsmath,amssymb}
\usepackage{minted}
\usepackage{multirow}
\usepackage{float}
\usepackage{algorithm}
\usepackage{algpseudocode}
\usepackage{todonotes}
\usepackage{forest}
\usepackage{xstring}
\usepackage{xcolor}
\usepackage[most]{tcolorbox}
\usepackage{xparse}
\usepackage{natbib}

\title{Visual existing graphic iPad edit context}
\date{2026-09-13}
\author{Mohammad Rahmani}

\begin{document}
\maketitle

\section{Problem}
Right-clicking an already rendered image in the Visual LaTeX editor could show
\texttt{Insert image in iPad} instead of \texttt{Edit image in iPad}. The Visual
context menu intentionally identifies an existing image only when the clicked
block carries the canonical \texttt{\_BlockMeta} metadata type with
\texttt. The context-normalization controller previously changed
only the \texttt{kind} attribute of whatever user-data object happened to be on
the block. If that object was absent or was a compatible non-canonical metadata
object, the menu's type check still rejected it and displayed the insertion
label.

\section{Fix}
\begin{itemize}
   \item Existing Visual image blocks are normalized to the exact \texttt{\_BlockMeta} type expected by the context-menu contract before the menu is built.
   \item Existing canonical metadata is retained in place and only normalized to \texttt.
   \item Compatible legacy/non-canonical metadata is copied into a new canonical graphic metadata object so source-line and source-segment information are preserved.
   \item If the source line is missing, the controller recovers it from the Visual source-position mapping so the existing backend \texttt{find\_reference} lookup still opens the correct image for editing.
   \item Detection works both from an \texttt{\textbackslash includegraphics} raw block and from the QText image character format.
   \item The underlying Insert/Edit request handler and iPad bridge workflow are unchanged.
\end{itemize}

\section{Regression coverage}
A regression test verifies that the controller normalizes image metadata to the
canonical type, preserves/reconstructs the source location needed by the edit
lookup, and leaves the existing Visual menu contract for both Edit and Insert in
place.

\section{Repository hygiene}
The correction stays inside the existing Visual graphic context controller,
which already owns this normalization responsibility. No runtime patch or
parallel image workflow was introduced.

\end{document}
